\chapter{Testowanie hipotez: konstrukcja, moc i interpretacja}
\label{ch:testing}
\paragraph{Powiązania i odsyłacze.} Teoria decyzji pojawia się w rozdziale~\ref{ch:decision}; asymptotyczne narzędzia testowania — w rozdziale~\ref{ch:asymptotic-estimation}.


\section{Od pytania praktycznego do testu}
Zakład wykazuje w ostatnim kwartale $7$ wad na $100$ losowo wybranych
produktów. Standard jakości wymaga, by prawdopodobieństwo wady nie
przekraczało $0{,}1$. Czy zaobserwowany odsetek $0{,}07$ świadczy o
poprawie? Trzeba odróżnić trzy pytania: (i) czy prawdopodobieństwo
wady jest mniejsze od $0{,}1$; (ii) czy jest ono \emph{znacząco}
mniejsze w sensie statystycznym; (iii) czy różnica jest wystarczająco
duża, by uzasadnić zmianę procesu produkcyjnego. Trzecie pytanie
wymaga także kryterium decyzji i kosztów, omówionych w rozdziale~\ref{ch:decision}.

Test jest regułą zależną od próby, a nie prawdopodobieństwem
prawdziwości hipotezy. Aby test miał sens, musimy zadeklarować
model, kierunek alternatywy, poziom istotności i sposób zebrania danych
\emph{przed} oglądaniem wyniku, jeśli chcemy zachować deklarowaną
kontrolę błędów.

\begin{definicja}[Hipotezy i test randomizowany]\label{def:statistical-test}
Niech $\mathcal E=(\mathsf X,\mathcal A,\{P_\theta:\theta\in\Theta\})$
będzie eksperymentem statystycznym, a $\Theta_0,\Theta_1$
niepustymi, rozłącznymi podzbiorami $\Theta$.
Hipotezę $H_0:\theta\in\Theta_0$ nazywamy zerową, a
$H_1:\theta\in\Theta_1$ alternatywną. \emph{Testem
randomizowanym} jest mierzalna funkcja
$\varphi:\mathsf X\to[0,1]$: po obserwacji $x$ hipotezę $H_0$
odrzucamy z prawdopodobieństwem $\varphi(x)$, korzystając w razie
potrzeby z niezależnego losowania pomocniczego. Test jest
\emph{nierandomizowany}, jeśli przyjmuje tylko wartości $0$ i $1$.
Funkcja $\beta_\varphi(\theta):=\E_\theta\varphi(X)$ jest
\emph{funkcją mocy}, natomiast
$\sup_{\theta\in\Theta_0}\beta_\varphi(\theta)$ jest
\emph{rozmiarem testu}. Test ma \emph{poziom istotności} co najwyżej
$\alpha\in[0,1]$, jeżeli $\sup_{\theta\in\Theta_0}\beta_\varphi(\theta)
\le\alpha$.
\end{definicja}

Jeśli $\theta\in\Theta_0$, odrzucenie jest błędem I rodzaju; jego
prawdopodobieństwo jest kontrolowane przez poziom testu. Jeśli
$\theta\in\Theta_1$, nieodrzucenie jest błędem II rodzaju;
prawdopodobieństwo tego błędu wynosi $1-\beta_\varphi(\theta)$.
Sama wartość $\alpha$ nie mówi nic o mocy ani o rozmiarze efektu.
Nieodrzucenie $H_0$ nie oznacza jego zaakceptowania jako prawdziwego.

\begin{przyklad}[Trzy pytania, trzy hipotezy]
Dla $X_i\iid\operatorname{Bernoulli}(p)$ można rozważać
$H_0:p\ge 0{,}1$ przeciw $H_1:p<0{,}1$ (poprawa),
$H_0:p\le 0{,}1$ przeciw $H_1:p>0{,}1$ (pogorszenie),
lub $H_0:p=0{,}1$ przeciw $H_1:p\ne0{,}1$ (dowolna zmiana).
To są różne problemy; wynik dwustronnego testu nie zastępuje
jednostronnego testu zaplanowanego ze względu na ryzyko wad.
\end{przyklad}

\section{Lemat Neymana--Pearsona: skąd bierze się najlepsza reguła?}
\begin{twierdzenie}[Lemat Neymana--Pearsona, z randomizacją]
\label{thm:neyman-pearson}
Niech $P_0,P_1$ będą dwoma miarami probabilistycznymi na
$(\mathsf X,\mathcal A)$ i niech $P_1\ll P_0$. Oznaczmy
$L=\dd P_1/\dd P_0$ ich iloraz wiarygodności. Niech
$0<\alpha<1$. Istnieją $c\ge0$ i $\gamma\in[0,1]$ takie, że
\[
 \varphi_*(x)=\ind{\{L(x)>c\}}+
 \gamma\ind{\{L(x)=c\}},\qquad \E_0\varphi_*=\alpha.
\]
Wtedy dla każdego mierzalnego testu $\varphi:\mathsf X\to[0,1]$
spełniającego $\E_0\varphi\le\alpha$ zachodzi
$\E_1\varphi\le\E_1\varphi_*$. A zatem $\varphi_*$ ma
największą moc wśród testów poziomu co najwyżej $\alpha$ dla
hipotez prostych $H_0:P=P_0$ i $H_1:P=P_1$.
\end{twierdzenie}
\begin{proof}
Ponieważ $L\ge0$ i $\E_0L=1$, istnieje próg $c\ge0$ spełniający
$P_0(L>c)\le\alpha\le P_0(L\ge c)$. Można go dobrać
jako odpowiedni kwantyl rozkładu $L$ pod $P_0$; gdy $c=0$,
warunek ma sens, bo $P_0(L\ge0)=1$. Jeśli
$P_0(L=c)>0$, ustalamy
$\gamma=[\alpha-P_0(L>c)]/P_0(L=c)$, a w przeciwnym razie
$\gamma=0$. Nierówności zapewniają $\gamma\in[0,1]$ i
$\E_0\varphi_*=\alpha$.

W punktach, w których $L>c$, mamy $\varphi_*=1$, zatem
$(L-c)(\varphi_*-\varphi)\ge0$. Przy $L<c$ mamy
$\varphi_*=0$, co znów daje nieujemny iloczyn; na zbiorze
$L=c$ iloczyn znika. W rezultacie
\[
 \E_1(\varphi_*-\varphi)
 =\E_0[L(\varphi_*-\varphi)]
 \ge c\E_0(\varphi_*-\varphi)\ge0.
\]
Wykorzystaliśmy absolutną ciągłość, by zapisać wartość
oczekiwaną pod $P_1$ jako całkę względem $P_0$, oraz kontrolę
poziomu testu, a nie szczególną postać obserwacji.
\end{proof}

\begin{uwaga}[Dlaczego randomizacja?]
Dla dyskretnych rozkładów nie zawsze można znaleźć zbiór
odrzucenia o prawdopodobieństwie dokładnie $\alpha$ pod $H_0$.
Test bez randomizacji zwykle ma \emph{rzeczywisty} rozmiar
mniejszy niż deklarowany poziom; jest to poprawny test
konserwatywny. Randomizacja na granicy pozwala wykorzystać
cały dopuszczony budżet błędu I rodzaju, lecz w typowych analizach
praktycznych częściej wybiera się prostszą procedurę
nierandomizowaną.
\end{uwaga}

\section{Testy jednostajnie najmocniejsze i statystyki monotoniczne}
\begin{definicja}[Test jednostajnie najmocniejszy]
Test $\varphi_*$ o poziomie co najwyżej $\alpha$ jest
\emph{jednostajnie najmocniejszy} (UMP) dla
$H_0:\theta\in\Theta_0$ przeciw $H_1:\theta\in\Theta_1$,
jeśli dla każdego $\theta\in\Theta_1$ i każdego testu
$\varphi$ tego samego poziomu mamy
$\E_\theta\varphi\le\E_\theta\varphi_*$.
\end{definicja}

\begin{twierdzenie}[Jednostronny test dwumianowy jest UMP]
\label{thm:binomial-ump}
Niech $X_1,\ldots,X_n\iid\operatorname{Bernoulli}(p)$,
$0<p<1$, $S=\sum_iX_i$, a $p_0\in(0,1)$.
Dla $0<\alpha<1$ wybierzmy $k\in\{0,\ldots,n\}$ oraz
$\gamma\in[0,1]$ tak, żeby funkcja
$\varphi_*(X)=\ind{\{S>k\}}+\gamma\ind{\{S=k\}}$
spełniała $\E_{p_0}\varphi_*=\alpha$.
Wtedy $\varphi_*$ jest testem UMP poziomu $\alpha$ dla
$H_0:p\le p_0$ przeciw $H_1:p>p_0$.
\end{twierdzenie}
\begin{proof}
Dla dowolnego $p_1>p_0$ iloraz wiarygodności próby wynosi
\[
 L(x)=\left(\frac{p_1}{p_0}\right)^{S(x)}
 \left(\frac{1-p_1}{1-p_0}\right)^{n-S(x)}.
\]
Stosunek $L(S+1)/L(S)=p_1(1-p_0)/[p_0(1-p_1)]>1$,
więc test Neymana--Pearsona przeciw $p_1$ odrzuca dla
dużych wartości $S$ i ma graniczną randomizację.
Ponieważ rozkład $S$ przy $p_0$ nie zależy od $p_1$,
można wybrać ten sam próg $k$ i $\gamma$ dla wszystkich
$p_1>p_0$.

Trzeba jeszcze sprawdzić poziom względem całego $H_0$.
Niech $U_1,\ldots,U_n$ będą niezależne z rozkładu
$U(0,1)$ i zdefiniujmy $S_p=\sum_i\ind{\{U_i\le p\}}$.
Dla $p\le p_0$ zachodzi $S_p\le S_{p_0}$ prawie na pewno.
Ponieważ funkcja $s\mapsto\ind{\{s>k\}}+
\gamma\ind{\{s=k\}}$ jest niemalejąca, moc przy $p$
nie przewyższa mocy przy $p_0$, czyli $\alpha$.
Lemat Neymana--Pearsona zapewnia następnie największą moc
przy każdym $p_1>p_0$ wśród testów poziomu $\alpha$.
\end{proof}

\begin{uwaga}[Test dwustronny]
W modelu normalnym o znanej wariancji test obustronny
$H_0:\mu=\mu_0$ przeciw $H_1:\mu\ne\mu_0$
nie ma na ogół testu UMP: najbardziej mocna reguła przeciw
$\mu_1>\mu_0$ odrzuca dla dużej średniej, natomiast przeciw
$\mu_1<\mu_0$ dla małej średniej. Wspólny test o maksymalnej
mocy dla obu kierunków nie może jednocześnie być obydwoma
jednostronnymi testami Neymana--Pearsona. Potrzebne są dodatkowe
kryteria, np. nieobciążoność testu lub niezmienniczość.
\end{uwaga}

\section{Wartość p: definicja i częste błędy interpretacji}
\begin{definicja}[Wartość p dla statystyki testowej]\label{def:pvalue}
Niech $T:\mathsf X\to\R$ będzie mierzalną statystyką,
a większe jej wartości świadczą przeciw $H_0$.
Dla obserwacji $x$ definiujemy
\[
 p(x):=\sup_{\theta\in\Theta_0}
 P_\theta\bigl(T(X)\ge T(x)\bigr).
\]
Zakładamy, że funkcja $x\mapsto p(x)$ jest mierzalna.
Dla hipotezy prostej supremum jest niepotrzebne.
\end{definicja}

\begin{twierdzenie}[Kontrola błędu przez wartość p]
\label{thm:pvalue-valid}
Jeśli $p$ jest określona jak w definicji~\ref{def:pvalue},
to dla każdego $u\in[0,1]$ oraz każdego $\theta\in\Theta_0$
zachodzi $P_\theta(p(X)\le u)\le u$.
\end{twierdzenie}
\begin{proof}
Ustalmy $\theta\in\Theta_0$ i oznaczmy
$G(t)=P_\theta(T(X)\ge t)$.
Z definicji $p(X)\ge G(T(X))$ prawie na pewno, a więc
$\{p(X)\le u\}\subseteq\{G(T(X))\le u\}$.
Dla dowolnej rzeczywistej zmiennej losowej $Y$ funkcja
$G_Y(t)=P(Y\ge t)$ jest nierosnąca. Zbiór
$A_u=\{t\in\R:G_Y(t)\le u\}$ jest zatem pusty,
cały lub jest półprostą postaci $[a,\infty)$ bądź
$(a,\infty)$. W pierwszym przypadku
$P(Y\in A_u)=0$. W drugim, o ile $A_u=[a,\infty)$,
mamy $P(Y\in A_u)=G_Y(a)\le u$.
Jeśli $A_u=(a,\infty)$, to z ciągłości prawdopodobieństwa
przy granicach monotonicznych
$P(Y>a)=\lim_{m\to\infty}P(Y\ge a+1/m)\le u$.
Dla $A_u=\R$ zachodzi $u=1$.
Zatem zawsze $P(G_Y(Y)\le u)\le u$. Stosując to do
$Y=T(X)$, otrzymujemy
$P_\theta(G(T(X))\le u)\le u$, a więc tezę.
\end{proof}

Dla ciągłej statystyki o ciągłym rozkładzie pod prostym
$H_0$ wartość p jest jednostajna na $[0,1]$; dla statystyk
dyskretnych najczęściej jest \emph{nadjednostajna} w sensie
powyższej nierówności, zatem dokładna nominalna równość nie
musi zachodzić. W żadnym przypadku $p=0{,}03$ nie oznacza
$3\%$ prawdopodobieństwa, że hipoteza zerowa jest prawdziwa.

\begin{uwaga}[Jedno- i dwustronne wartości p nie są tym samym]
Definicja~\ref{def:pvalue} jest jednostronnym ogonem statystyki,
której duże wartości przemawiają przeciw hipotezie zerowej.
W symetrycznym ciągłym modelu normalnym dla
$H_0:\mu=\mu_0$ i statystyki $Z$ o rozkładzie
$\Normal(0,1)$ dwustronna wartość p ma postać
$P(|Z'|\ge|z_{\rm obs}|)=2[1-\Phi(|z_{\rm obs}|)]$,
gdzie $Z'\sim\Normal(0,1)$. W modelu dyskretnym,
np. dwumianowym, istnieje kilka konwencji dotyczących
„tak samo ekstremalnych” wyników dwustronnych, które
nie muszą dawać identycznych wartości p. Należy więc
zawsze nazwać zastosowaną regułę odrzucenia, a nie tylko
podać słowo „dwustronny”.
\end{uwaga}

\begin{zadanie}[7 wad na 100 produktów: test dokładny]\label{zad:binom-test}
W modelu $X_i\iid\operatorname{Bernoulli}(p)$
zaobserwowano $S=7$ przy $n=100$. Testujemy
$H_0:p\ge0{,}1$ przeciw $H_1:p<0{,}1$.
Uzasadnij postać wartości p i oblicz ją.
\par\emph{Rozwiązanie.} Mniejsza liczba wad przemawia za
alternatywą, zatem z monotoniczności rozkładu dwumianowego
w parametrze $p$ mamy
\[
 p\text{-value}=\sup_{p\ge0{,}1}P_p(S\le7)
 =P_{0{,}1}(S\le7)
 =\sum_{j=0}^7\binom{100}{j}(0{,}1)^j(0{,}9)^{100-j}
 \approx0{,}2061.
\]
Na poziomie $5\%$ nie odrzucamy $H_0$. Nie oznacza to,
że udowodniliśmy $p\ge0{,}1$: badanie może mieć
niewystarczającą moc dla niewielkiej poprawy.
\end{zadanie}

\section{Moc, wielkość efektu i dobór liczebności próby}
\begin{twierdzenie}[Dokładna moc jednostronnego testu normalnego]
\label{thm:normal-power}
Niech $X_1,\ldots,X_n\iid\Normal(\mu,\sigma^2)$,
$\sigma>0$ będzie znana, a $\mu_0\in\R$.
Test odrzucający $H_0:\mu\le\mu_0$ na rzecz
$H_1:\mu>\mu_0$ wtedy i tylko wtedy, gdy
$\bar X>\mu_0+z_{1-\alpha}\sigma/\sqrt n$, ma
poziom $\alpha$ i funkcję mocy
\[
 \beta(\mu)=1-\Phi\left(z_{1-\alpha}
 -\frac{\sqrt n(\mu-\mu_0)}{\sigma}\right),
\]
gdzie $\Phi$ jest dystrybuantą rozkładu $\Normal(0,1)$.
\end{twierdzenie}
\begin{proof}
Z dokładnego rozkładu $\bar X\sim\Normal(\mu,\sigma^2/n)$
po standaryzacji otrzymujemy podaną funkcję mocy.
Jest ona rosnąca w $\mu$, a więc jej największa wartość na
$\{\mu\le\mu_0\}$ występuje w $\mu_0$ i wynosi $\alpha$.
\end{proof}

\begin{wniosek}[Dobór liczebności dla zadanego efektu]
\label{cor:normal-sample-size}
Przy założeniach twierdzenia~\ref{thm:normal-power}, dla
$0<\alpha<1/2$, $0<\beta_0<1/2$ i efektu $\delta>0$
warunek mocy co najmniej $1-\beta_0$ przy $\mu=\mu_0+\delta$
jest równoważny
\[
 n\ge\left(\frac{\sigma(z_{1-\alpha}+z_{1-\beta_0})}
 {\delta}\right)^2.
\]
Najmniejsze dopuszczalne $n$ uzyskujemy przez
zaokrąglenie prawej strony w górę (z zastrzeżeniem $n\ge1$).
\end{wniosek}
\begin{proof}
Nierówność $1-\Phi(z_{1-\alpha}-\sqrt n\delta/\sigma)
\ge1-\beta_0$ oznacza
$z_{1-\alpha}-\sqrt n\delta/\sigma\le z_{\beta_0}
=-z_{1-\beta_0}$. Po przestawieniu otrzymujemy tezę.
\end{proof}

\begin{zadanie}[Ile pomiarów potrzeba?]
Załóżmy znane odchylenie standardowe $\sigma=10$ jednostek.
Chcemy z poziomem istotności $\alpha=0{,}05$ wykryć wzrost
średniej o $\delta=3$ z mocą co najmniej $0{,}80$.
\par\emph{Rozwiązanie.} Dla $z_{0{,}95}\approx1{,}645$
oraz $z_{0{,}80}\approx0{,}842$ wystarcza
$n\ge[10(1{,}645+0{,}842)/3]^2\approx68{,}73$,
czyli $n=69$. To rachunek dla \emph{znanej} $\sigma$,
niezależnych obserwacji i testu jednostronnego. Przy danych
zależnych, większym rzeczywistym rozproszeniu lub kilku
porównaniach potrzebny jest inny projekt.
\end{zadanie}

Nie należy określać $\delta$ jako „istotnego efektu” wyłącznie
dlatego, że dla dostatecznie dużego $n$ daje małą wartość p.
Trzeba wcześniej ustalić \emph{minimalny efekt praktycznie ważny},
czyli najmniejszą zmianę, która miałaby znaczenie dla decyzji.
Raportowanie przedziału ufności i miary wielkości efektu jest
zwykle znacznie bardziej informacyjne niż samo odrzucenie hipotezy.

\section{Pułapki praktyczne: selekcja, wielokrotność i zatrzymywanie}
\begin{przyklad}[Zatrzymywanie, gdy pojawi się mała wartość p]
Załóżmy, że pięć niezależnych badań prowadzonych pod prawdziwą
hipotezą zerową daje niezależne wartości p o rozkładzie
jednostajnym na $[0,1]$. Jeśli ogłosimy sukces, gdy
\emph{którakolwiek} wartość będzie mniejsza od $0{,}05$,
prawdopodobieństwo przynajmniej jednego fałszywego alarmu wynosi
$1-(1-0{,}05)^5\approx0{,}2262$, a nie $0{,}05$.
Ten sam problem pojawia się przy wielokrotnym testowaniu albo
ciągłym podglądaniu danych i zakończeniu eksperymentu w
pierwszej korzystnej chwili. Zależności między testami mogą
zmienić tę liczbę, ale nie usprawiedliwiają ignorowania problemu.
\end{przyklad}

\begin{center}\small
\begin{tabular}{>{\raggedright\arraybackslash}p{0.30\textwidth}>{\raggedright\arraybackslash}p{0.57\textwidth}}
\toprule
\textbf{Pytanie przed testem} & \textbf{Dlaczego jest ważne?}\\
\midrule
Czy hipotezę i kierunek wybrano wcześniej? & Wybranie kierunku po zobaczeniu danych może naruszać kontrolę poziomu.\\
Czy obserwacje są niezależne? & Korelacja wewnątrz grup zwykle zmienia rozkład statystyki i moc.\\
Czy testujemy różnicę, przewagę czy równoważność? & Brak istotnej różnicy nie jest dowodem równoważności.\\
Czy rozkład statystyki jest dokładny czy asymptotyczny? & Małe próby i rzadkie zdarzenia mogą wymagać testu dokładnego.\\
Ile hipotez sprawdzamy? & Wielokrotność wymaga jawnej strategii kontroli błędów.\\
Czy efekt jest użyteczny? & Wartość p nie mierzy wielkości ani praktycznej ważności efektu.\\
\bottomrule
\end{tabular}
\end{center}

\begin{lstlisting}[language=Python,caption={Dokładny test dwumianowy i obliczenie mocy testu normalnego.}]
from math import ceil, sqrt
from scipy.stats import binomtest, norm

k, n, p0 = 7, 100, 0.10
result = binomtest(k, n, p0, alternative="less")
print("p-value:", result.pvalue)

alpha, power, sigma, delta = 0.05, 0.80, 10.0, 3.0
z_alpha = norm.ppf(1 - alpha)
z_power = norm.ppf(power)
n_required = ceil((sigma * (z_alpha + z_power) / delta)**2)
print("minimalne n:", n_required)
\end{lstlisting}

\paragraph{Co sprawdzamy przed uruchomieniem programu?}
Pakiet \texttt{scipy} należy zainstalować w środowisku Pythona;
kod zakłada niezależne próby Bernoulliego oraz znaną wariancję
w planowaniu testu normalnego. Żadne wywołanie funkcji testowej
nie sprawdza automatycznie reprezentatywności próby ani
prawidłowości projektu badania.


% ================================================================
% ROZDZIAL 10 -- TESTY PARAMETRYCZNE I WPROWADZENIE DO ANOVA
% ================================================================
